\section{Arithmetic Fourier charts and coherent tower refinement}
\label{sec:arithmetic-fourier-charts}

The inclusion of a finite field and the uniform pullback along its trace
give two quantum codes in the larger register. Their exact relative angle
depends on whether the extension degree vanishes in the characteristic.
All degrees admit a uniform Gram classification. A positive conditional
matrix chart survives when its angle closes, and degrees invertible in the
characteristic admit coherent refinement through towers of every finite
length. The three positive results have passed the capped review; the full
mixed arithmetic construction remains conjectural.
The checker is \path{theory/checks/arithmetic_limits_check.py}.

\subsection{The exact finite-field sectors}

\begin{definition}[Arithmetic register and transfer conventions]
Fix a prime $p$ and a named primitive complex $p$th root $\zeta_p$.
For a finite field $E/\mathbb F_p$ of cardinality $p^r$, retain the field,
the Hilbert space $H_E=\ell^2(E)$ with orthonormal basis $|x\rangle$,
and the character
\[
 \psi_E(x)=\zeta_p^{\operatorname{Tr}_{E/\mathbb F_p}(x)},\qquad
 \operatorname{Tr}_{E/\mathbb F_p}(x)=\sum_{j=0}^{r-1}x^{p^j}.
\]
The prime-field phase space is $E\oplus E$ with alternating form
$\Omega_E((a,b),(a',b'))=\operatorname{Tr}_{E/\mathbb F_p}(ab'-a'b)$.
The reference operators are
\[
 X_E(a)|x\rangle=|x+a\rangle,\quad
 Z_E(b)|x\rangle=\psi_E(bx)|x\rangle,\quad
 W_E(a,b)=Z_E(-b)X_E(a).
\]
An ordered register list has tensor-product Hilbert space and direct-sum
phase space; the empty list has Hilbert space $\mathbb C$. Field structures
and factor order remain part of the data. Frobenius and Fourier are
\[
 U_E|x\rangle=|x^p\rangle,\qquad
 F_E|x\rangle=|E|^{-1/2}\sum_{y\in E}\psi_E(-xy)|y\rangle.
\]
Frobenius on a list is their tensor product. A subscript on a Fourier
tensor factor means identity on the other factors.
For a named embedding $i:K\to E$ over the same $p,\zeta_p$, with
$|K|=p^s$, $|E|=p^{ns}$, put
\[
 T_i(y)=i^{-1}\!\left(\sum_{j=0}^{n-1}y^{p^{sj}}\right),\qquad
 \kappa_i=|\ker T_i|,
\]
where the inverse is restricted to $i(K)$. Composite traces use their
composite named embeddings. The inclusion and trace-fibre transfers are
\[
 J_i|a\rangle=|i(a)\rangle,\qquad
 V_i|a\rangle=\kappa_i^{-1/2}\sum_{T_i(y)=a}|y\rangle.
\]
All square roots are positive, and reverse arrows are Hilbert adjoints.
\end{definition}
\begin{scope}
The negative Fourier kernel and the common named root of unity are fixed.
These are finite-field registers; a separable embedding-set register over
an arbitrary field is a separate Galois construction.
\end{scope}
\provenance{D1301--D1304, D1306}{foundation.md; transfers-example.md}
{arithmetic Frobenius checker}{frobenius-hierarchy-adjudication.md}

\begin{definition}[Joint arithmetic code and Gram datum]
For a named embedding $i:K\to E$ of degree $n\geq2$, put
$q=|K|=p^s$, $Q=|E|=q^n$ and $\kappa_i=Q/q$. Write
\[
 \Gamma_i=J_i^*V_i,\quad P_i=J_iJ_i^*,\quad
 \widetilde P_i=V_iV_i^*,\quad
 \mathcal S_i=\operatorname{ran}J_i+\operatorname{ran}V_i.
\]
Let $E_i$ project orthogonally onto $\mathcal S_i$, and use the physical
trace $\tau_E=\operatorname{Tr}/Q$. The projection algebra is the unital
star algebra generated by $P_i,\widetilde P_i$ in $\operatorname{End}(H_E)$;
zero-dimensional summands are omitted. Put
$|+_K\rangle=q^{-1/2}\sum_x|x\rangle$ and $R_K|x\rangle=|-x\rangle$.
The joint Frobenius representation means $U_E$ restricted to $\mathcal S_i$.
\end{definition}
\begin{scope}
The joint range is a subspace of the larger register, with its inherited
physical trace. Its normalized conditional trace will be specified separately.
\end{scope}
\provenance{D1421}{definitions.md}{arithmetic limits checker F1--F2}{arithmetic-mixed-adjudication.md}

\begin{definition}[Corrected trace frame and singular sectors]
If $p\nmid n$, define
\begin{gather*}
 D_n^K|x\rangle=|nx\rangle,\qquad L_i=V_iD_n^K,\qquad
 c_i=\kappa_i^{-1/2},\qquad h_i=\sqrt{1-c_i^2},\\
 W_i=\frac{L_i-c_iJ_i}{h_i},\qquad B_i=(D_n^K)^{-1}F_K.
\end{gather*}
The map $\mathcal T_i:\mathbb C^2\otimes H_K\to H_E$ has columns
$J_i,W_i$. For any nonzero prime-field scalar $a$, scalar dilations on
each field are $D_a|x\rangle=|ax\rangle$.
If $p\mid n$, instead put
\[
 \alpha_i=J_i|+_K\rangle,\quad\beta_i=V_i|0\rangle,\quad
 d_i=\sqrt{q/\kappa_i},\quad A_K=|+_K\rangle^\perp.
\]
For $n>2$ put $e_i=\sqrt{1-d_i^2}$ and let the angle frame
$\mathcal T_i^{\rm ang}:\mathbb C^2\to H_E$ have columns
$\alpha_i,(\beta_i-d_i\alpha_i)/e_i$. For every $p\mid n$ the residual
frame $\mathcal T_i^{\rm res}:\mathbb C^2\otimes A_K\to H_E$ has columns
$J_i|_{A_K},V_iF_K|_{A_K}$. If $p=2,n=2$, retain the common line
$\mathbb C\alpha_i=\mathbb C\beta_i$ and only the residual frame.
\end{definition}
\begin{scope}
The corrected degree dilation is defined only when the degree is invertible.
The singular quadratic case has no angle frame requiring division by zero.
\end{scope}
\provenance{D1422, D1423}{definitions.md}{arithmetic limits checker F1--F2}{arithmetic-mixed-adjudication.md}

\begin{proposition}[Uniform Gram sectors, physical weights and symmetry]
\statusProved\quad
For every named finite-field embedding $i:K\to E$ of degree $n\geq2$,
\[
 \sqrt{\kappa_i}(\Gamma_i)_{x,a}=\delta_{a,nx}.
\]
If $p\nmid n$, all $q$ singular-value squares are $1/\kappa_i$.
If $p\mid n$, there is one square $q/\kappa_i=q^{2-n}$ and $q-1$ zeros.
The ranges intersect only when $p=2,n=2$, in the invariant line
$\mathbb C\alpha_i$.
The joint projection algebra and its physical block weights are
\[
\begin{array}{c|c|c}
 \text{degree stratum}&\text{joint algebra, with multiplicities}&
 \text{physical weights}\\ \hline
 p\nmid n&M_2\otimes I_q&2q/Q\\
 p\mid n,\ n>2&M_2\oplus\mathbb C\oplus\mathbb C\ (1,q-1,q-1)&
 2/Q,\ (q-1)/Q,\ (q-1)/Q\\
 p=2,\ n=2&\mathbb C\oplus\mathbb C\oplus\mathbb C\ (1,q-1,q-1)&
 1/Q,\ (q-1)/Q,\ (q-1)/Q.
\end{array}
\]
Each matrix summand uses its normalized matrix trace. The complement is
one scalar summand, of weight $1-2q/Q$ in the first two cases and
$1-(2q-1)/Q$ in the last.
In the regular frame Fourier and Frobenius are
\[
 \mathcal T_i^*F_E\mathcal T_i
 =\begin{pmatrix}c_i&h_i\\h_i&-c_i\end{pmatrix}\otimes B_i,
 \qquad
 \mathcal T_i^*U_E\mathcal T_i=I_2\otimes U_K.
\]
On the singular angle plane Fourier is the same matrix with parameters
$d_i,e_i$, while Frobenius is identity. On the singular residual frame,
\[
 F_{\rm res}=\begin{pmatrix}0&R_K|_{A_K}\\I_{A_K}&0\end{pmatrix},
 \qquad U_{\rm res}=I_2\otimes U_K|_{A_K}.
\]
In quadratic characteristic two the common line is Fourier-fixed and
$F_{\rm res}=X\otimes I_{A_K}$, where
$X=\left(\begin{smallmatrix}0&1\\1&0\end{smallmatrix}\right)$.
The joint Frobenius representation is two copies of $U_K$, with one
invariant common line removed in that exceptional case. In every stratum
$U_E^s=I$ on $\mathcal S_i$ when $q=p^s$.
\end{proposition}
\begin{scope}
These formulas hold for all finite fields and named embeddings of degree
at least two. The projection algebra alone is commutative in the exceptional
case; its residual matrix factor additionally uses compressed Fourier.
The complement retains the actual ambient Fourier and Frobenius restrictions.
\end{scope}
\provenance{MIX-GRAM}{gram-and-chart.md, Sections 1--4}
{arithmetic limits checker F1--F2}{arithmetic-mixed-adjudication.md}

\begin{proof}[Derivation]
The trace of $i(x)$ is $nx$, so evaluating a normalized trace-fibre column
at $i(x)$ gives the Gram entry. If $n$ is invertible it is a scaled
permutation matrix. Otherwise it is
$\sqrt{q/\kappa_i}|+_K\rangle\langle0|$. A common range vector corresponds
exactly to singular value one, giving the stated exception.
Gram subtraction produces the displayed orthogonal frames. The negative
Fourier kernel gives $F_EJ_i=V_iF_K$, $F_EV_i=J_iF_K$ and $F_K^2=R_K$;
also $F_KD_n^K=(D_n^K)^{-1}F_K$. Substitution gives every Fourier block,
including the residual reflection in odd characteristic. Frobenius
intertwines both transfers and fixes the uniform and zero vectors.
The block weights are their dimensions divided by $Q$.
\end{proof}

\subsection{A positive conditional chart and its tangent}

\begin{definition}[Conditioned matrix chart]
For $0\leq c\leq1$ let $h=\sqrt{1-c^2}$, use $M_2(\mathbb C)$ with
faithful reference $\operatorname{tr}_2=\operatorname{Tr}/2$, and set
\[
 p_c=\begin{pmatrix}1&0\\0&0\end{pmatrix},\quad
 q_c=\begin{pmatrix}c^2&ch\\ch&h^2\end{pmatrix},\quad
 f_c=\begin{pmatrix}c&h\\h&-c\end{pmatrix},\quad
 d_c=\begin{pmatrix}-h&c\\c&h\end{pmatrix}.
\]
For $c<1$, $d_c=(q_c-p_c)/h$; at one retain its continuous displayed
value $d_1=X$, while $f_1=Z=\operatorname{diag}(1,-1)$.
Conditioning means dividing a represented chart's restricted physical
trace by its support weight before varying $c$. With a retained logical
space $H$, use $\operatorname{tr}_2\otimes\operatorname{tr}_H$ on
$M_2\otimes\operatorname{End}(H)$, with chart matrices on the first factor.
\end{definition}
\begin{scope}
This is a matrix continuation in an angle parameter. No field of
noninteger cardinality is posited.
\end{scope}
\provenance{D1424}{definitions.md}{arithmetic limits checker R1--R2}{arithmetic-mixed-adjudication.md}

\begin{definition}[Continuous instruments and retained decoding]
\leavevmode\par
For finite-dimensional Hilbert spaces $X,Y_b$, a continuous instrument
on a closed real parameter interval is a finite family of continuous
matrices $K_b(t):X\to Y_b$ satisfying $\sum_bK_b(t)^*K_b(t)=I_X$.
It sends $\rho$ to $\bigoplus_bK_b(t)\rho K_b(t)^*$, with ordinary matrix
traces on the output blocks. For a continuous isometry $S(t):X\to Y$,
the retained decoder amplitudes are
\[ S(t)^*:Y\to X,\qquad I_Y-S(t)S(t)^*:Y\to Y, \]
with distinct success and failure tags, respectively.
Independent instruments use tensor-product Kraus amplitudes and retain all
outcome tuples. Sequential instruments retain every history. A reference
state supplies input probabilities and does not replace the ordinary
trace in the channel normalization equation.
\end{definition}
\begin{scope}
Failure retains the ambient state. Distinct sequential and direct
measurement protocols are not identified by forgetting their histories.
\end{scope}
\provenance{D1426}{definitions.md}{arithmetic limits checker R2--R3}{arithmetic-mixed-adjudication.md}

\begin{proposition}[Faithful Fourier tangent after conditioning]
\statusProved\par
For every regular-degree joint chart and every singular-degree angle plane
with $n>2$, physical conditioning gives the matrix chart just defined,
with its stated logical factor when present. Throughout the closed angle
interval,
\[
 d_c^2=f_c^2=I,\qquad f_cp_cf_c=q_c,\qquad f_cd_c=-d_cf_c.
\]
At $c=1$ the two projections coincide, while the retained tangent $X$
and $p_1$ generate $M_2$ with faithful trace $\operatorname{Tr}/2$.
Actual chart Fourier becomes $Z$ and conjugates $X$ to $-X$.
With the regular logical factor retained, the full Fourier operator is
$Z\otimes B_i$ at this endpoint.
Continuous instruments, retained decoders and independent tensor remain
completely positive and trace preserving at the endpoint.
Quadratic characteristic two has instead a common invariant line and a
residual factor $M_2\otimes I_{A_K}$ generated by its compressed projections
and Fourier $X\otimes I_{A_K}$, with normalized residual reference
$\operatorname{tr}_2\otimes\operatorname{tr}_{A_K}$.
\end{proposition}
\begin{scope}
The coalescing-plane continuation has the declared regular or singular
nonquadratic scope. The exceptional residual factor is positive but is
not obtained by dividing by a zero angle. Frobenius remains on the
logical factor, and is identity on the singular angle plane.
\end{scope}
\provenance{MIX-CHART}{gram-and-chart.md, Sections 5--6}
{arithmetic limits checker F2, R1--R2}{arithmetic-mixed-adjudication.md}

\begin{proof}[Derivation]
The four matrix identities follow by multiplication. The trace is faithful
because $\operatorname{tr}_2(A^*A)=\frac12\sum_{j,k}|A_{jk}|^2$.
At one, preparing $p_1$ and testing $p_1$ gives probability one, whereas
applying $d_1=X$ first gives zero. Complete positivity follows because
every amplification is a sum of positive quadratic forms
$(I\otimes K_b)\rho(I\otimes K_b)^*$; the Kraus completeness equation
gives trace preservation, including for tensor products.
For the rank-one columns $e_0$ and $v=ce_0+he_1$, $f_ce_0=v$ gives
the full Fourier decoder comparison on both success and failure.
\end{proof}

The order of operations matters. The formal unconditioned complement
weight in the regular branch is $1-2/\kappa$, which is $-1/3$ at
$\kappa=3/2$. The positive continuation above first conditions the
arithmetic chart and then varies its angle; it does not continue that
whole physical trace as a positive reference state.

\subsection{Refinement through arbitrary finite towers}

\begin{definition}[Path frame and refinement connector]
Let $K_0\xrightarrow{i_1}\cdots\xrightarrow{i_\ell}K_\ell$, $\ell\geq1$,
be a named tower with degrees $n_r\geq2$ invertible in characteristic $p$.
Set $N=\prod_rn_r$. In outer-to-inner chart order $\ell,\ldots,1$, let
\[
 \mathcal T_{\rm path}=\mathcal T_{i_\ell}
 (I_2\otimes\mathcal T_{i_{\ell-1}})\cdots
 (I_{2^{\ell-1}}\otimes\mathcal T_{i_1})
 : (\mathbb C^2)^{\otimes\ell}\otimes H_{K_0}\longrightarrow H_{K_\ell}.
\]
For any $0<c_r<1$ put $h_r=\sqrt{1-c_r^2}$,
$v_r=c_re_0+h_re_1$, $c_*=\prod_rc_r$, $h_*=\sqrt{1-c_*^2}$.
Define $C:\mathbb C^2\to(\mathbb C^2)^{\otimes\ell}$ by
\[
 Ce_0=|0\cdots0\rangle,\qquad
 Ce_1=\frac{v_\ell\otimes\cdots\otimes v_1-c_*|0\cdots0\rangle}{h_*}.
\]
At arithmetic fibres use $c_r=\kappa_{i_r}^{-1/2}$. Consecutive blocks
use the products of their constituent parameters in the same order.
\end{definition}
\begin{scope}
Only the path frame requires a physical tower. The connector is also a
well-defined matrix family at nonarithmetic parameter values.
\end{scope}
\provenance{D1425}{definitions.md}{arithmetic limits checker F3, R3}{arithmetic-mixed-adjudication.md}

\begin{definition}[Weighted endpoint connector]
For fixed positive weights $a_1,\ldots,a_\ell$, put $A=\sum_ra_r$ and
$c_r(t)=t^{-a_r/2}$ for $t>1$. At $t=1$ define
\[
 C_a(1)e_0=|0\cdots0\rangle,\qquad
 C_a(1)e_1=\sum_r\sqrt{\frac{a_r}{A}}\,|\operatorname{one}_r\rangle,
\]
where $|\operatorname{one}_r\rangle$ has its sole entry one in edge r's
factor. Consecutive blocks have summed weights. The direct chart uses
$\operatorname{tr}_2$; the whole path space uses
$\operatorname{tr}_{2^\ell}$. Conditioning the latter on the connector
range gives the former. Their unconditioned reference states are distinct.
\end{definition}
\begin{scope}
Every weight is positive. Identity edges may be omitted; they require
no local chart factor.
\end{scope}
\provenance{D1427}{definitions.md}{arithmetic limits checker R3}{arithmetic-mixed-adjudication.md}

\begin{proposition}[Coherent refinement and Fourier through every tower length]
\statusProved\quad
For every finite named tower of length $\ell\geq1$ and degrees
$n_r\geq2$ invertible in the characteristic, the corrected trace maps
compose exactly. The path
frame and connector are isometries. For composite embedding k,
\[
 L_{i_\ell}\cdots L_{i_1}=L_k,\qquad
 \mathcal T_k=\mathcal T_{\rm path}(C\otimes I_{H_{K_0}}).
\]
All consecutive-block refinements give the same connector into the common
path basis. The actual Fourier and Frobenius restrictions satisfy
\begin{align*}
 F_{K_\ell}\mathcal T_{\rm path}
 &=\mathcal T_{\rm path}\left[
 (f_{c_\ell}\otimes\cdots\otimes f_{c_1})
 \otimes (D_N^{K_0})^{-1}F_{K_0}\right],\\
 U_{K_\ell}\mathcal T_{\rm path}
 &=\mathcal T_{\rm path}(I_{2^\ell}\otimes U_{K_0}),\\
 (f_{c_\ell}\otimes\cdots\otimes f_{c_1})C&=Cf_{c_*}.
\end{align*}
For every positive weight vector the connector extends continuously to
the displayed weighted endpoint, and refinement and Fourier identities
continue there; in particular $Z^{\otimes\ell}C_a(1)=C_a(1)Z$.
Retained decoder success amplitudes compose exactly. Full sequential
history instruments and all independent tensor outcomes stay continuous,
completely positive and trace preserving through the endpoint.
\end{proposition}
\begin{scope}
There is no bound on the finite tower length. Singular-degree refinements
are not supplied by this formula. The direct connector occupies a
two-dimensional subspace of the full independent path tensor. Equality of
successful decoding does not identify a sequential history instrument with
the binary decoder of the composite isometry.
\end{scope}
\provenance{MIX-TOWER}{coherent-towers.md, Sections 1--6}
{arithmetic limits checker F3, R3}{arithmetic-mixed-adjudication.md}

\begin{proof}[Derivation]
Trace linearity gives $D_a^LV_i=V_iD_a^K$ for every nonzero prime-field
scalar. Thus for consecutive degrees m,n,
$V_jD_n^LV_iD_m^K=V_jV_iD_{mn}^K=L_{ji}$.
The path frame sends the empty path to the direct inclusion and
$v_\ell\otimes\cdots\otimes v_1$ to the direct corrected trace column;
Gram subtraction gives C. Each block connector preserves these two
vectors, so every grouping gives the same linear map. Fourier exchanges
the two vectors in each factor and scalar dilations commute through the
local frames, giving the precise logical scaling shown above.
Writing $t=1+\epsilon$ gives
$1-t^{-a_r}=a_r\epsilon+O(\epsilon^2)$. The coefficient of a path
with k excitations is of order $\epsilon^{(k-1)/2}$ in the second column.
Only single excitations survive, with the stated square-root weights.
Finally, an isometry S has retained Kraus completeness
$SS^*+(I-SS^*)^2=I$; all endpoint and tensor assertions follow continuously.
\end{proof}

For two edges j after i the nontrivial connector column is
\[
 \frac{c_jh_i|01\rangle+h_jc_i|10\rangle+h_jh_i|11\rangle}
 {\sqrt{1-c_j^2c_i^2}}.
\]
For three weights $(1,2,4)$ the endpoint is
$\sqrt{1/7}|\operatorname{one}_1\rangle+
\sqrt{2/7}|\operatorname{one}_2\rangle+
\sqrt{4/7}|\operatorname{one}_3\rangle$.
The encoded tangent has square $CC^*$, rather than identity on the entire
path tensor. Higher excitation sectors remain available for independent
processes, even though direct refinement does not fill them.

\subsection{The remaining mixed arithmetic problem}

\begin{definition}[Finite mixed extension test diagram]
Take a finite diagram of named finite-field embeddings in one
characteristic, closed under its named composites, with finite ordered
register words. Its named amplitude operations are Frobenius, Fourier,
inclusion and trace transfer and their adjoints, nonzero prime-field
dilations, tensor identities, and the reversible multiplication gates
\[
 M_E^{(d)}|x_1,\ldots,x_d,z\rangle
 =|x_1,\ldots,x_d,z+\textstyle\prod_jx_j\rangle
\]
for each explicitly named positive finite arity $d\geq1$, with their finite composites and
adjoints. Retain their isometric encodings and decoding instruments,
every sequential outcome history and every independent outcome tuple.
Marked sector tests include the joint and singular sectors just defined,
explicitly selected Frobenius orbit projections and active-control
projections, and their actual Fourier transports. An orbit projection
selects its named Frobenius basis orbits; an active-control projection
requires each named control to be nonzero and leaves the target arbitrary.
An equality is equality of the named arithmetic amplitude maps or tagged
CP maps, not equality created by compressing intermediate gates separately.
\end{definition}
\begin{scope}
Measurement placement is part of the source process. Retaining measurement
histories does not identify a measured circuit with an unmeasured circuit.
\end{scope}
\provenance{D1428, D1308, D1327}{definitions.md}
{arithmetic limits checker M1}{arithmetic-mixed-adjudication.md}

\begin{conjecture}[A positive mixed arithmetic process realization]
\statusConjecture\quad
There is one positive filtered process realization natural in every finite
mixed extension test diagram, including all characteristic-dividing tower
sectors, and compatible with the separately specified general Galois
embedding-set fragment through its stated orbit comparisons. It retains
the named Fourier-transfer and multiplication relations, all instrument
histories and independent tensor outcomes, and recovers the prescribed
conditional chart, orbit and active-sector normalizations.
Such a realization must specify its operator algebras, CP arrows, reference
weights, continuation rule for mixed moments, and diagram comparison maps.
It must give positive mixed moment Gram matrices at every finite matrix
amplification and exhibit retained Frobenius and active multiplication
state/effect distinctions in the common model.
\end{conjecture}
\begin{scope}
No such full realization is constructed here. The all-characteristic Gram
classification and invertible-degree tower refinement supply constraints
on it, while an all-characteristic tower connector and a compatible
mixed-moment normalization remain open. Separate constant matrix envelopes
or unrelated local continuations do not establish the conjecture.
\end{scope}
\provenance{MIX-ALL}{coherent-towers.md, Section 7}
{arithmetic limits checker M1 is only a necessary finite-fibre diagnostic}{arithmetic-mixed-adjudication.md}

The next explicit comparison uses both
$\mathbb F_2\to\mathbb F_4\to\mathbb F_{16}$ and
$\mathbb F_3\to\mathbb F_9\to\mathbb F_{81}$, adding
$\mathbb F_3\to\mathbb F_{27}$ for its singular odd-characteristic
reflection. The arithmetic matrices already distinguish the full
multiplication circuit from the circuit obtained by deleting every
intermediate sector exit. A candidate must retain those exits and their
subsequent transitions while preserving the probability of each specified
measurement protocol. That is the next mixed requirement beyond the
positive tower chart established here.
